\chapter{Laminectomy}
\index{Laminectomy}

\section*{Definition and Overview}
Laminectomy is a major decompressive surgical procedure that involves the surgical removal of the lamina --- the flat, bony arch forming the posterior roof of the spinal canal --- from one or more vertebrae. The primary clinical objective of this procedure is to create additional space within the spinal canal, thereby relieving pathological pressure on the spinal cord, cauda equina, or spinal nerve roots. It is most commonly performed in the lumbar or cervical spine to treat spinal stenosis, herniated discs, spinal tumours, or bony overgrowth (osteophytes). By decompressing these neural structures, the procedure aims to alleviate debilitating pain, sensory abnormalities, and motor weakness, restoring physical function and improving the patient's quality of life.

\section*{Epidemiology}
Degenerative spinal diseases, particularly lumbar spinal stenosis, represent the leading indication for laminectomy, with a peak incidence in individuals aged 60 years and older. In many countries and globally, the volume of spinal decompression surgeries has steadily risen due to an ageing population and increased clinical detection via advanced imaging. While laminectomy is performed across all age groups, its utilisation is highly correlated with age-related degenerative changes. The procedure is slightly more prevalent in males than females, reflecting a higher rate of symptomatic spinal stenosis in the male demographic.

\section*{Aetiology and Pathophysiology}
Laminectomy is indicated when spinal canal narrowing compromises neural structures, resulting in clinical symptoms that fail to respond to conservative therapy. The pathophysiology involves mechanical compression and vascular compromise of the spinal cord or nerve roots. Indicated causes and risk factors include:

\begin{itemize}
    \item \textbf{Lumbar spinal stenosis:} Age-related hypertrophic changes of the facet joints, ligamentum flavum, and intervertebral discs that narrow the central spinal canal and lateral recesses.
    \item \textbf{Herniated intervertebral disc:} Extrusion of the nucleus pulposus through the annulus fibrosus, causing focal compression on adjacent nerve roots.
    \item \textbf{Spondylolisthesis:} The forward displacement of one vertebra over another, causing mechanical narrowing of the spinal canal.
    \item \textbf{Spinal tumours:} Benign or malignant neoplasms within or adjacent to the spinal canal that exert compressive force on the spinal cord.
    \item \textbf{Trauma:} Vertebral fractures, dislocations, or acute disc herniations resulting from high-impact injuries.
\end{itemize}

\section*{Clinical Features}
The presentation of patients requiring laminectomy varies depending on the spinal level affected, but typically reflects progressive nerve or spinal cord compression. The typical clinical signs and symptoms include:

\begin{itemize}
    \item \textbf{Neurogenic claudication:} Leg pain, numbness, or heaviness that is exacerbated by walking or standing and relieved by sitting or bending forward (the 'shopping trolley sign').
    \item \textbf{Radiculopathy:} Sharp, radiating pain, parasthesia, or sensory loss along a specific dermatomal distribution in the lower limbs (lumbar) or upper limbs (cervical).
    \item \textbf{Motor weakness:} Weakness in specific muscle groups, such as foot drop (associated with L5 nerve root compression) or reduced hand grip strength (cervical).
    \item \textbf{Reflex changes:} Diminished deep tendon reflexes at the knee or ankle, or pathologically brisk reflexes with clonus if spinal cord compression is present (myelopathy).
    \item \textbf{Postoperative features:} Temporary localised pain, stiffness, and wound tenderness at the surgical site are normal clinical features during the early recovery phase.
\end{itemize}

\section*{Differential Diagnosis}
Before proceeding to surgery, it is critical to distinguish spinal stenosis from other conditions that mimic its clinical presentation. These include:

\begin{itemize}
    \item \textbf{Vascular claudication:} Lower limb pain brought on by exertion, but caused by peripheral arterial disease rather than spinal compression; it is relieved rapidly by standing still and is associated with weak peripheral pulses.
    \item \textbf{Hip or knee osteoarthritis:} Degenerative joint disease causing groin, thigh, or knee pain that can resemble radicular spinal pain but is aggravated by joint rotation and weight-bearing.
    \item \textbf{Peripheral neuropathy:} Distal, symmetrical sensory loss and pain (often in a 'glove-and-stocking' distribution), commonly caused by diabetes mellitus or vitamin deficiencies.
    \item \textbf{Trochanteric bursitis:} Localised inflammation over the greater trochanter, causing lateral hip pain that is tender to palpation and does not radiate below the knee.
    \item \textbf{Meningioma or schwannoma:} Intradural, extramedullary spinal tumours that may present with slow, progressive radicular pain or myelopathy, requiring differentiation from simple degenerative stenosis.
\end{itemize}

\section*{When to Seek Medical Care}
Patients undergoing conservative management for spinal stenosis, or those recovering postoperatively, must monitor for signs of progressive neurological deterioration or surgical complications.

\textbf{Contact your local emergency services for:}
\begin{itemize}
    \item Sudden loss of bowel or bladder control, or saddle anaesthesia (numbness in the perineum, buttocks, and inner thighs), indicating cauda equina syndrome.
    \item Rapidly progressive, severe weakness in one or both legs or arms.
    \item Clear, watery fluid leaking from the surgical wound, suggesting a cerebrospinal fluid leak.
    \item Signs of a deep vein thrombosis or pulmonary embolism, such as sudden shortness of breath, chest pain, or severe swelling and calf tenderness.
    \item High fever, worsening neck stiffness, or severe headache.
\end{itemize}

\section*{Diagnosis and Investigations}
A structured diagnostic workup is essential to confirm anatomical stenosis and correlate it with the patient's clinical symptoms. The investigations include:

\begin{itemize}
    \item \textbf{Magnetic resonance imaging (MRI):} The gold standard diagnostic modality, offering detailed images of soft tissues, including spinal cord compression, ligamentum flavum hypertrophy, and disc herniation.
    \item \textbf{Computed tomography (CT) scan:} Provides detailed bone anatomy, useful for assessing osteophyte formation, facet joint arthropathy, and in patients where MRI is contraindicated.
    \item \textbf{CT myelography:} An invasive imaging study where contrast dye is injected into the spinal canal prior to a CT scan, highly effective for defining spinal canal narrowing.
    \item \textbf{Plain spinal X-rays:} Utilised to evaluate overall spinal alignment, degenerative disc disease, and dynamic instability via flexion and extension views.
    \item \textbf{Electromyography (EMG) and Nerve Conduction Studies (NCS):} Performed to assess nerve function and differentiate spinal nerve root compression from peripheral nerve entrapment or neuropathy.
\end{itemize}

\section*{Management}
Management is initially conservative, but decompressive surgery is indicated if there is progressive neurological deficit or severe pain unresponsive to non-operative measures. Strategies include:

\begin{itemize}
    \item \textbf{Surgical decompression (laminectomy):} An open or minimally invasive procedure performed under general anaesthesia, involving the resection of the lamina, thickened ligamentum flavum, and compressive osteophytes.
    \item \textbf{Pharmacological therapy:} Postoperative pain is managed using multimodal analgesia, including paracetamol, non-steroidal anti-inflammatory drugs (NSAIDs), and short-term oral opioids.
    \item \textbf{Physiotherapy and rehabilitation:} Graded postoperative exercise programmes focusing on core stabilisation, lumbar mobility, and general conditioning to restore function.
    \item \textbf{Epidural corticosteroid injections:} A non-surgical management option to provide temporary pain relief by reducing inflammation around compressed nerve roots.
    \item \textbf{Lifestyle and activity modification:} Gradual return to activities, avoiding heavy lifting, prolonged sitting, or lumbar twisting during the immediate recovery phase.
\end{itemize}

\section*{Complications}
While laminectomy is generally safe, it carries risks inherent to spinal surgery. Potential complications include:

\begin{itemize}
    \item \textbf{Dural tear:} Accidental puncture of the dura mater during surgery, leading to a cerebrospinal fluid (CSF) leak, which may cause postural headaches and require surgical repair.
    \item \textbf{Nerve root damage:} Permanent or temporary injury to the spinal nerves, resulting in new or worsened motor weakness, sensory deficit, or chronic neuropathic pain.
    \item \textbf{Surgical site infection:} Infection of the superficial skin or deep spinal tissues, presenting with localised redness, warmth, and purulent discharge, requiring antibiotics or wound debridement.
    \item \textbf{Spinal instability:} Loss of structural support from excessive bone removal, which may necessitate a future spinal fusion procedure.
    \item \textbf{Failed back surgery syndrome:} Persistent back or leg pain despite adequate surgical decompression of the spinal canal.
    \item \textbf{Venous thromboembolism:} Deep vein thrombosis (DVT) or pulmonary embolism (PE) secondary to perioperative immobility.
\end{itemize}

\section*{Prognosis and Outcomes}
The prognosis following a laminectomy is generally favourable, with approximately 70\% to 80\% of patients experiencing significant relief from leg pain and neurogenic claudication. Complete recovery of muscle strength and sensory function depends on the severity and duration of the preoperative nerve compression. Most patients are discharged from the hospital within one to three days and can gradually resume light activities within four to six weeks. Long-term outcomes can be influenced by the progression of degenerative disc and facet disease at adjacent spinal segments.

\section*{Prevention and Self-Care}
Patients can take several steps to protect their spine and support long-term surgical success:

\begin{itemize}
    \item \textbf{Maintain a healthy body weight:} Minimises the mechanical load and stress placed on the lumbar spine.
    \item \textbf{Regular low-impact exercise:} Engaging in walking, swimming, or stationary cycling to maintain joint mobility and cardiovascular fitness.
    \item \textbf{Core muscle strengthening:} Performing exercises that target the abdominal and back muscles to provide intrinsic spinal support.
    \item \textbf{Ergonomic lift training:} Utilising correct body mechanics, such as bending at the knees rather than the waist when lifting objects.
    \item \textbf{Smoking cessation:} Nicotine restricts microvascular blood flow, which impairs bone fusion, wound healing, and disc health.
\end{itemize}

\section*{Sources and Further Reading}
The following reputable organisations provide further information on laminectomy and spinal health:

\begin{itemize}
    \item Healthdirect Australia. \textit{Information on spinal decompression surgery, preparation, and recovery.} https://www.healthdirect.gov.au/
    \item Mayo Clinic. \textit{Clinical overview of laminectomy indications, risks, and post-surgical expectations.} https://www.mayoclinic.org/
    \item NHS. \textit{Patient guide to lumbar decompression surgery, including what happens during the procedure.} https://www.nhs.uk/
    \item Cleveland Clinic. \textit{Comprehensive details on surgical techniques and post-operative rehabilitation for laminectomy.} https://my.clevelandclinic.org/
    \item American Association of Neurological Surgeons. \textit{Patient education resources regarding lumbar spinal stenosis and surgical options.} https://www.aans.org/
\end{itemize}
